%%% FIELD clinical-statement
Liberali et al. \cite{LiberaliEtAl2025RTAR} report retrospective, empirically grounded resampling studies of GUSTO-1 and EUROPA: patients are drawn with replacement from original randomized-trial arm pools under a stationarity/exchangeability benchmark, and the displayed uncertainty intervals are the empirical \(2.5\%\) and \(97.5\%\) quantiles across \(200\) replicates.  GUSTO-1 contributes three analyzed arms and a death-or-survival endpoint at \(30\) days; EUROPA contributes two arms and a composite cardiovascular endpoint, with a reported \(4.2\)-year mean time to the last observation.  Their baseline real-time rule makes daily Gittins-index assignments, whereas their eta-variant mixes in equal randomization only for arms that have not reached fixed minimum counts.  Those published percentile intervals are neither \(I_{T,\Delta}^{\mathrm W}\) nor \(I_{T,\Delta}\).  Consequently, this note implies no coverage or tail-padding claim for the reported intervals.  The published allocation description certifies neither the permanent floors \(\pi_t(1)\geq ct^{-\alpha}\) and \(\pi_t(0)\geq\varepsilon\) nor a known polynomial envelope with parameters \((\beta,L)\); the reported EUROPA mean time to the last observation is not such a tail envelope.  Hence no numerical EUROPA tail pad follows.  For GUSTO-1 only, if complete \(30\)-day ascertainment and this note's single-reveal encoding hold, observation for at least \(30\) days after the final randomization makes \(D_t(a)\leq30\) for every realized assigned endpoint and therefore makes \(b_{T,\Delta}=0\) exactly; that fact alone does not establish adaptive-Wald validity.  Finally, consider instead a prospective two-arm Liberali-style experiment modified to maintain the two permanent logged floors, target the risk difference \(\tau=P Y_t(1)-P Y_t(0)\), use this note's interval, and certify the stated independent-arrival, bounded-outcome, sequential-randomization, and polynomial-tail conditions.  Then \(I_{T,\Delta}\) is uniformly asymptotically honest over \(\mathfrak M_T^{\mathrm{low}}(\alpha,\beta,L,c,\varepsilon)\); on the nonempty regular subclass \(\mathfrak M_{T,\Delta}^{\mathrm{reg}}(\kappa)\), ordinary Wald is uniformly asymptotically honest if and only if
\[
 T^q(1+T+\Delta_T)^{-\beta}\longrightarrow0,
 \qquad q=(1-\alpha)/2.
\]
Thus no extra follow-up is needed when \(\beta>q\); when \(\beta=q\), one needs \(\Delta_T/T\to\infty\); and when \(\beta<q\), one needs \(T+\Delta_T\gg T^{q/\beta}\).  Transfer to an odds-ratio target or to the three-arm GUSTO-1 design is outside this proposition.

%%% FIELD theorem-statement
For fixed \(\eta\in(0,1/2)\), fixed class parameters \(\alpha\in[0,1)\), \(\beta,L,c,\varepsilon>0\) satisfying the declared parameter restrictions, and any prespecified deterministic \(\Delta_T\geq0\), there exist constants \(0<c_\eta\leq C_\eta<\infty\) such that
\[
 \inf_{(P,\Pi,T)\in\mathfrak M_T^{\mathrm{low}}(\alpha,\beta,L,c,\varepsilon)}
 P\{\tau\in I_{T,\Delta}\}\geq1-\eta-o(1)
\]
and
\[
 \sup_{(P,\Pi,T)\in\mathfrak M_T^{\mathrm{low}}(\alpha,\beta,L,c,\varepsilon)}
 \mathbb E_{P,\Pi}|I_{T,\Delta}|\leq C_\eta r_{T,\Delta}.
\]
Conversely, every \(O_{T,\Delta}\)-measurable interval \(J_{T,\Delta}\) satisfying
\[
 \inf_{(P,\Pi,T)\in\mathfrak M_T^{\mathrm{low}}(\alpha,\beta,L,c,\varepsilon)}
 P\{\tau\in J_{T,\Delta}\}\geq1-\eta-o(1)
\]
obeys
\[
 \sup_{(P,\Pi,T)\in\mathfrak M_T^{\mathrm{low}}(\alpha,\beta,L,c,\varepsilon)}
 \mathbb E_{P,\Pi}|J_{T,\Delta}|\geq c_\eta r_{T,\Delta}
\]
for all sufficiently large \(T\). Here \(r_{T,\Delta}=\max\{T^{-(1-\alpha)/2},(1+T+\Delta_T)^{-\beta}\}\), and each expectation integrates over the independent and identically distributed outcome--delay arrivals under \(P\) and every randomization induced by the predictable policy \(\Pi\). Both converse components use the same fixed policy \(\pi_t(1)=ct^{-\alpha}\): prompt mean tilts give the sampling term, while finite hidden-sign delays at \(H_T+1\) give the tail term. Thus, at \(\Delta_T=0\), the exact class-minimax expected-length order is \(\max\{T^{-(1-\alpha)/2},T^{-\beta}\}\).

%%% FIELD proof-clinical
The source facts are recorded in \(\texttt{prop:liberali-main-source-record}\) and \(\texttt{prop:liberali-supplement-source-record}\), transcribed from the main manuscript and online supplement.  In particular, the experiment in that paper is a retrospective resampling experiment: conditional on a chosen arm, a record is sampled with replacement from the corresponding original-trial arm pool, the procedure is repeated \(200\) times, and the reported lower and upper limits are empirical replicate quantiles.  By contrast, \(\texttt{def:ordinary-wald}\) defines, for one realized logged experiment,
\[
 I_{T,\Delta}^{\mathrm W}
 =\bigl[\widehat\tau_{T,\Delta}
 \mathbin{\pm}z_{1-\eta/2}\widehat s_{T,\Delta}\bigr]\cap[-2,2],
 \tag{1}
\]
and \(\texttt{def:hybrid-interval}\) defines \(I_{T,\Delta}\) by an observable variance gate, a bounded-shift Gaussian radius containing \(B_{T,\Delta}\), and a concentration fallback.  An empirical quantile of \(200\) resampling replicates is not either map in (1) or \(\texttt{def:hybrid-interval}\).  The coverage theorems for those two defined maps therefore make no assertion about the percentile intervals reported by Liberali et al.

The same source records show why the model assumptions cannot be imputed to those simulations.  The baseline rule assigns the patients arriving on a day according to a Gittins-index choice.  Its description does not provide a positive lower bound for every arm on every day.  In the eta-variant, the equal-randomization component is applied only while an arm remains below a predetermined minimum count; after all minima have been attained, that description supplies no permanent exploration or control floor.  Thus neither
\[
 \pi_t(1)\geq ct^{-\alpha}
 \quad\text{nor}\quad
 \pi_t(0)\geq\varepsilon
 \tag{2}
\]
is certified by the published rule.  Likewise, the reported \(4.2\)-year EUROPA quantity is an empirical mean calendar time to the last observation, not an armwise population statement of the form
\[
 P\{D_t(a)>u\}\leq\min\{1,L(1+u)^{-\beta}\}
 \quad\text{for every integer }u\geq0.
 \tag{3}
\]
It therefore identifies neither \(L\) nor \(\beta\), and the numerical pad \(B_{T,\Delta}\) cannot be evaluated from that mean.

For the limited GUSTO-1 conclusion, impose the two conditions stated in the proposition: complete ascertainment of the binary \(30\)-day endpoint and the single-reveal encoding.  Waiting at least \(30\) days after final randomization means \(\Delta_T\geq30\).  Since \(k_T\leq T\),
\[
 h_{T,\Delta}=T+\Delta_T-k_T\geq30.
 \tag{4}
\]
Complete ascertainment gives \(D_t(a)\leq30\) for every realized assigned endpoint, so (4) and \(\texttt{def:bias-radius}\) imply
\[
 \begin{aligned}
 b_{T,\Delta}
 &=P\bigl[Y(1)\mathbf 1\{D(1)>h_{T,\Delta}\}\bigr]
   -P\bigl[Y(0)\mathbf 1\{D(0)>h_{T,\Delta}\}\bigr]\\
 &=0.
 \end{aligned}
 \tag{5}
\]
Equation (5) removes delayed-tail bias, but it supplies none of the overlap, variance, or adaptive-limit conditions needed for ordinary Wald, so it cannot by itself establish Wald coverage.

It remains to prove the conditional prospective transfer.  Under the modifications enumerated in the proposition, the independent and identically distributed potential outcome--delay vectors satisfy \(\texttt{ass:iid-arrivals}\), binary outcomes satisfy \(\texttt{ass:bounded-outcomes}\), the known logs satisfy \(\texttt{ass:sequential-randomization}\), (2) is imposed as \(\texttt{ass:one-sided-forced-exploration}\) and \(\texttt{ass:control-floor}\), and (3) is imposed as \(\texttt{ass:polynomial-delay-envelope}\).  Hence the experiment belongs to \(\mathfrak M_T^{\mathrm{low}}(\alpha,\beta,L,c,\varepsilon)\) by \(\texttt{def:one-sided-policy-class}\), and for binary outcomes its causal risk difference is
\[
 P\{Y_t(1)=1\}-P\{Y_t(0)=1\}
 =P Y_t(1)-P Y_t(0)=\tau.
 \tag{6}
\]
The uniform honesty of the hybrid interval now follows directly from \(\texttt{thm:uniform-gated-inference}\), including its quota-failure convention and concentration fallback.  If the regular-variance condition also holds and the regular class is nonempty, \(\texttt{prop:wald-followup-frontier}\) gives the necessary and sufficient ordinary-Wald condition
\[
 x_T:=T^q(1+T+\Delta_T)^{-\beta}\longrightarrow0,
 \qquad q=(1-\alpha)/2>0.
 \tag{7}
\]
The three regimes are algebraic consequences of (7).  If \(\beta>q\), choosing \(\Delta_T=0\) gives \(x_T=T^{q-\beta}\{1+o(1)\}\to0\).  If \(\beta=q\),
\[
 x_T=\left(\frac{T}{1+T+\Delta_T}\right)^q\to0
 \quad\Longleftrightarrow\quad
 \frac{\Delta_T}{T}\to\infty.
 \tag{8}
\]
If \(\beta<q\), positivity permits taking the \(\beta\)-th root, and
\[
 x_T\to0
 \quad\Longleftrightarrow\quad
 \frac{1+T+\Delta_T}{T^{q/\beta}}\to\infty
 \quad\Longleftrightarrow\quad
 T+\Delta_T\gg T^{q/\beta}.
 \tag{9}
\]
All deductions in (6)--(9) concern the prospective two-arm risk-difference protocol after the stated modifications.  Neither the odds-ratio estimands reported by Liberali et al. nor the three-arm GUSTO-1 design is identified with this proposition's target.

%%% FIELD proof-theorem
Write \(q=(1-\alpha)/2\).  First note the design-class relation.  The two-sided class is contained in the one-sided class because its rare-arm lower bound is the same and
\[
 \pi_t(0)=1-\pi_t(1)\geq1-Ct^{-\alpha}\geq1-C\geq\varepsilon.
 \tag{1}
\]
In particular, the deterministic policy
\[
 p_t:=\pi_t(1)=ct^{-\alpha},
 \qquad \pi_t(0)=1-p_t,
 \tag{2}
\]
belongs to \(\mathfrak M_T^{\mathrm{low}}(\alpha,\beta,L,c,\varepsilon)\), since \(1-p_t\geq1-c\geq\varepsilon\).

For the upper bound, uniform coverage is the one-sided-class conclusion of \(\texttt{thm:uniform-gated-inference}\).  By \(\texttt{lem:quota-success}\), uniformly over that class,
\[
 K_1(n_T)\asymp T^{1-\alpha},\quad K_0(n_T)\asymp T,
 \quad J_1\asymp T^{\rho(1-\alpha)},\quad J_0\asymp T^\rho,
 \quad P(E_T^c)\leq4T^{-4}.
 \tag{3}
\]
Consequently,
\[
 S_T=K_1(n_T)^{-1}+K_0(n_T)^{-1}=O(T^{-(1-\alpha)}),
 \qquad a_T=O(T^{-q}).
 \tag{4}
\]
Because \(g_{T,\Delta}=m_T+\Delta_T\geq m_T\asymp T^\rho\), \(\texttt{ass:polynomial-delay-envelope}\) gives
\[
 v_{T,\Delta}=O(T^{-\rho\beta}).
 \tag{5}
\]
For every fixed \(x>0\), (3)--(5) and \(\texttt{def:hybrid-interval}\) imply
\[
 A_{1,T,\Delta}(x)
 =O\!\left(T^{-\rho(\beta+1-\alpha)/2}
          +T^{-\rho(1-\alpha)}\right)=o(T^{-q}),
 \tag{6}
\]
and
\[
 A_{0,T,\Delta}(x)
 =O\!\left(T^{-\rho(\beta+1)/2}+T^{-\rho}\right)=o(T^{-q}).
 \tag{7}
\]
Indeed, the chosen cohort exponent satisfies
\[
 \rho(\beta+1-\alpha)>1-\alpha,
 \qquad \rho(1-\alpha)>(1-\alpha)/2,
 \tag{8}
\]
and the arm-zero inequalities are weaker.  The strict inequalities in (8) absorb the factors \(\sqrt{\log T}\) and \(\log T\) appearing when \(x=4\log T\); hence
\[
 d_{T,\Delta}=o(T^{-q}).
 \tag{9}
\]

For observations in \([-1,1]\), the usual unbiased sample variance based on at least two observations is at most two.  Therefore, on \(E_T\),
\[
 \widehat s_{T,\Delta}^2\leq2S_T,
 \qquad \widehat s_{T,\Delta}=O(T^{-q}).
 \tag{10}
\]
Let \(z=z_{1-\eta/2}\).  At \(c=u+z\), the defining left-hand side for the bounded-shift critical value is at least
\(2\Phi(z)-1=1-\eta\).  Monotonicity in \(c\) and uniqueness in \(\texttt{def:bounded-shift-critical-value}\) therefore give
\[
 c_\eta(u)\leq u+z,
 \qquad
 \widehat s_{T,\Delta}c_\eta(B_{T,\Delta}/\widehat s_{T,\Delta})
 \leq B_{T,\Delta}+z\widehat s_{T,\Delta}.
 \tag{11}
\]
Moreover, uniformly over deterministic \(\Delta_T\geq0\),
\[
 \frac{1+h_{T,\Delta}}{1+H_T}
 =1-\frac{k_T}{1+H_T}=1+o(1),
 \qquad
 B_{T,\Delta}=O((1+H_T)^{-\beta}).
 \tag{12}
\]
Equations (4)--(12) show that the length of either active branch on \(E_T\) is bounded by a constant multiple of
\[
 r_{T,\Delta}=\max\{T^{-q},(1+H_T)^{-\beta}\}.
 \tag{13}
\]
On \(E_T^c\), the interval is \([-2,2]\), and (3) yields
\[
 4P(E_T^c)\leq16T^{-4}=o(T^{-q})=o(r_{T,\Delta}).
 \tag{14}
\]
Taking expectations in (13)--(14) proves the asserted uniform upper bound using only the two allocation floors.

For the sampling converse, retain the single policy (2), set both delays to zero, and set \(Y_t(0)=0\).  Under \(P_+\) and \(P_-\), let \(Y_t(1)\in\{-1,1\}\) have means \(d_T\) and \(-d_T\), respectively, where
\[
 Q_T=\sum_{t=1}^T p_t,
 \qquad d_T=\lambda Q_T^{-1/2},
 \qquad 0<\lambda<\frac{1-2\eta}{2\sqrt2}.
 \tag{15}
\]
These prompt-delay laws satisfy all support and tail conditions, and integral comparison gives \(Q_T\asymp T^{1-\alpha}\), hence \(d_T\asymp T^{-q}\).  Their causal contrasts are \(\tau_+=d_T\) and \(\tau_-=-d_T\).  Conditional on the assignment at time \(t\), the one-step Hellinger affinity is one for arm zero and \(\sqrt{1-d_T^2}\) for arm one.  Independence over arrivals and the fixed policy therefore give the finite-sample identity
\[
 \operatorname{Aff}(P_+^O,P_-^O)
 =\prod_{t=1}^T\left[1-p_t\{1-\sqrt{1-d_T^2}\}\right]
 \geq1-Q_Td_T^2=1-\lambda^2.
 \tag{16}
\]
For completeness, if \(p\) and \(q\) are densities of two laws, Cauchy--Schwarz gives
\[
 \begin{aligned}
 \|P-Q\|_{\mathrm{TV}}
 &=\frac12\int|\sqrt p-\sqrt q|(\sqrt p+\sqrt q)\\
 &\leq\frac12
 \left\{\int(\sqrt p-\sqrt q)^2
             \int(\sqrt p+\sqrt q)^2\right\}^{1/2}
 =\sqrt{1-\operatorname{Aff}(P,Q)^2}.
 \end{aligned}
 \tag{17}
\]
Combining (16)--(17),
\[
 \|P_+^O-P_-^O\|_{\mathrm{TV}}\leq\sqrt2\lambda.
 \tag{18}
\]
Let \(J_{T,\Delta}\) be any interval satisfying the stated uniform-honesty premise.  Transfer the event \(\{\tau_-\in J_{T,\Delta}\}\) from \(P_-^O\) to \(P_+^O\) by (18), then apply the union bound under \(P_+^O\).  This yields
\[
 P_+^O\{\tau_+,\tau_-\in J_{T,\Delta}\}
 \geq1-2\eta-o(1)-\sqrt2\lambda>0
 \tag{19}
\]
for all sufficiently large \(T\).  On the event in (19), the interval length is at least \(2d_T\); hence for a constant \(c_{\eta,1}>0\),
\[
 \sup_{(P,\Pi,T)\in\mathfrak M_T^{\mathrm{low}}}
 \mathbb E_{P,\Pi}|J_{T,\Delta}|
 \geq c_{\eta,1}T^{-q}.
 \tag{20}
\]

For the hidden-tail converse, retain the same policy (2) and set
\[
 \delta_T=\min\{1/4,L(1+H_T)^{-\beta}\}.
 \tag{21}
\]
Independently over arrivals and arms, draw a Bernoulli indicator of probability \(\delta_T\).  On failure set \((Y_t(a),D_t(a))=(0,0)\).  On success set \(D_t(a)=H_T+1\); under \(P_+\), set the successful arm-one outcome to \(+1\) and the successful arm-zero outcome to \(-1\), and reverse both signs under \(P_-\).  Outcomes are bounded, delays are finite, and for every integer \(u\geq0\),
\[
 P_\pm\{D_t(a)>u\}
 =\begin{cases}
 \delta_T,&0\leq u\leq H_T,\\
 0,&u\geq H_T+1,
 \end{cases}
 \leq\min\{1,L(1+u)^{-\beta}\}.
 \tag{22}
\]
Thus both laws belong to the one-sided class.  Every nonzero atom is hidden until after \(H_T\); conditional on each of the \(2^T\) assignment sequences, the observed zero records and missingness indicators have the same law under the two signs.  Consequently the complete observed laws coincide, while
\[
 \tau_+=2\delta_T,
 \qquad \tau_-=-2\delta_T.
 \tag{23}
\]
Honesty under their common observed law gives
\[
 P\{\tau_+,\tau_-\in J_{T,\Delta}\}
 \geq1-2\eta-o(1).
 \tag{24}
\]
On this event the length is at least \(4\delta_T\).  Since \(H_T\geq T\to\infty\), the second term in (21) is eventually active, and for a constant \(c_{\eta,2}>0\),
\[
 \sup_{(P,\Pi,T)\in\mathfrak M_T^{\mathrm{low}}}
 \mathbb E_{P,\Pi}|J_{T,\Delta}|
 \geq c_{\eta,2}(1+H_T)^{-\beta}.
 \tag{25}
\]
Equations (16) and (22)--(24) are finite-sample calculations and do not condition on quota success.  In the smallest nontrivial assignment experiment, \(T=2\), (16) is the product over the four assignment sequences, while (22)--(24) hold separately on each sequence; this exact enumeration rules out a hidden conditioning or adaptive-stopping step in either converse.  Combining (20) and (25) proves
\[
 \sup_{\mathfrak M_T^{\mathrm{low}}}\mathbb E|J_{T,\Delta}|
 \geq\min\{c_{\eta,1},c_{\eta,2}\}
 \max\{T^{-q},(1+H_T)^{-\beta}\}.
 \tag{26}
\]
Together with the upper bound, (26) proves the exact one-sided class-minimax order.  At \(\Delta_T=0\), \(H_T=T\), so this order is \(\max\{T^{-(1-\alpha)/2},T^{-\beta}\}\).

We finally record the theorem-level comparison required for its interpretation.  The source statement \(\texttt{prop:shi-delayed-inference-scope}\) records that Shi, Wang, and Wu \cite{ShiWangWu2023DelayedInference} assume independent and identically distributed potential-outcome vectors independent of contemporaneous action and delay, stationary arm-specific delay laws, \((2+\delta)\)-moments, and adaptive-weight conditions; their Theorems 4.1--4.3 prove DAIPW consistency, asymptotic normality, and variance consistency.  Their Corollary 4.4 verifies those conditions under a one-sided allocation floor, positive prompt mass, and a polynomial finite-delay tail with exponent at least one half, with a nonzero-margin relaxation in Remark 4.5.  It does not state uniform honesty or a minimax expected-length theorem.  In the present class, boundedness is stronger than their moment condition, and \(\texttt{ass:polynomial-delay-envelope}\) excludes infinite delays because
\[
 P\{D_t(a)=\infty\}
 \leq P\{D_t(a)>u\}
 \leq L(1+u)^{-\beta}\longrightarrow0.
 \tag{27}
\]
Conversely, the present class permits arbitrary within-arm outcome--delay dependence and every \(\beta>0\), both beyond their global independence condition and their displayed Corollary 4.4 regime.  Thus neither full class contains the other.  Only the design-side floor is imported here: the prompt construction (15) lies in their outcome--delay-independence submodel, whereas the hidden-sign construction (21)--(23) deliberately makes the outcome sign depend on whether the delay is hidden and therefore does not.

%%% FIELD liberali-main-statement
The main manuscript of Liberali et al. reports retrospective, empirically grounded resampling of the GUSTO-1 and EUROPA arm pools.  Conditional on the arm selected by the simulated real-time rule, patients are drawn with replacement under the authors' stationarity and exchangeability benchmark; each study uses \(200\) replicates, and reported medians and intervals are the empirical \(2.5\%\), \(50\%\), and \(97.5\%\) levels across replicates.  The manuscript describes GUSTO-1 as a three-analyzed-arm study with a \(30\)-day all-cause-mortality endpoint, EUROPA as a two-arm study with a composite cardiovascular endpoint and a reported \(4.2\)-year mean time to the last outcome observation, and the baseline real-time allocation as a daily Gittins-index rule.

%%% FIELD liberali-supplement-statement
The online supplement to Liberali et al. specifies death or survival at \(30\) days after randomization as the GUSTO-1 endpoint and a composite of cardiovascular mortality, nonfatal myocardial infarction, and resuscitated cardiac arrest as the EUROPA endpoint.  It states that all reported manuscript results average \(200\) replicates, that confidence limits are empirical lower and upper \(2.5\%\) replicate levels, that the baseline daily rule selects an arm using a Gittins index, and that the eta-variant's equal-randomization component applies only until prespecified arm minimum counts are reached.

%%% FIELD shi-statement
Shi, Wang, and Wu assume independent and identically distributed potential-outcome vectors independent of contemporaneous action and delay, together with stationary arm-specific delay laws.  Their DAIPW consistency, asymptotic-normality, and variance-consistency theorems require adaptive-weight conditions and a finite \((2+\delta)\)-moment.  Their Corollary 4.4 verifies those conditions under a one-sided allocation floor, positive prompt mass, and a polynomial restriction on finite delays with tail exponent at least one half; their Remark 4.5 relaxes the last restriction to an allocation-exponent-plus-tail-exponent condition when the bandit margin is nonzero.  The paper proves neither uniform coverage over the present informative-delay class nor a minimax expected-length order.
